\section{A nonflat positive example}
The next construction shows that all polynomial tensor maps can be
isomorphisms even when the polynomial module is not flat. Its finite
tensor-balancing identity is analogous to the one in Lazard's
epimorphism example \cite[Section 110.36, Tag 06RH]{StacksLazardExample};
the all-arity polynomial calculation is given here.

\begin{theorem}\label{thm:nonflat-forest-plethory}
There is a proper inclusion $R\subset S$ of zero-dimensional local
$\mathbb F_2$-algebras for which every nilpotent has square zero and
$R$ is not Noetherian, such that
\[
          \Omega(R,S)^{\otimes_R n}\xrightarrow{\sim}\Omega_n(R,S)
                         \qquad(n\ge1).
\]
Both $S$ and $\Omega(R,S)$ are nonflat over $R$. The effective
coefficient ring is $S$. Thus flatness is not a necessary condition
for the existence of the unique compatible polynomial plethory.
\end{theorem}
\begin{proof}
Let the vertices $v=(j,w)$ range over countably many full binary
trees, where $j\ge1$ and $w$ is a finite binary word. Denote the
children by $v0,v1$, and define
\[
 S=\mathbb F_2[x_v,y_v:v]/(x_v^2,y_v^2,y_v-x_vy_{v0}y_{v1}:v).
\]
This presentation has an injective direct-limit model. At level
$n$ use the squarefree polynomial algebra on $x_v$ of depth below
$n$ and $y_v$ of depth $n$. Replace a boundary $y_v$ at the next
level by $x_vy_{v0}y_{v1}$. The new triples are pairwise disjoint
and disjoint from the retained variables. Distinct squarefree
monomials therefore have distinct nonzero images. Every finite
expression and every displayed relation occurs at a finite level,
so the direct limit is exactly $S$.

Give $x_v$ degree $-1$ and $y_v$ degree $1$, and take
$R=\bigoplus_{d\le0}S_d$. This is a unital subring, and it is proper
because each nonzero $y_v$ has degree $1$. The augmentation ideals
of both rings consist of square-zero elements: this follows from
the squarefree monomial model and the Frobenius identity. Every
element with constant term $1$ is its own inverse. Hence both
rings are local, with one prime and residue field $\mathbb F_2$.
Write $N$ for the maximal ideal of $R$.

Products of degree-zero monomials $x_vy_{v0}$ in distinct tree
components are nonzero. Consequently $N^k\ne0$ for every $k$.
A finitely generated ideal generated by square-zero elements is
nilpotent, so $N$ is not finitely generated and $R$ is not Noetherian.

The relations also make $R\to S$ epic. In $S\otimes_R S$,
\[
 y_v\otimes1
 =y_{v0}\otimes x_vy_{v1}
 =x_vy_{v0}\otimes y_{v1}
 =1\otimes y_v,
\]
because both $x_vy_{v0}$ and $x_vy_{v1}$ belong to $R$.
The analogous identity for $x_v$ is automatic. These generators
give $S\otimes_R S\simeq S$.

For every $k\ge1$ one has the exact colon identity
\[
                              (R:_S N^k)=R.
\]
Indeed, if $c\notin R$, its nonzero positive-degree part involves
only finitely many components. Multiply it by a product of $k$
degree-zero monomials from fresh components. This product belongs
to $N^k$ and cannot kill that positive-degree part: at a finite
level multiplication across disjoint components is the tensor
product of nonzero vectors over $\mathbb F_2$, and the transition
maps are injective. The product remains outside $R$.

Set $Q(X)=X^4+X^2=X^2(X-1)^2$. For $r=a+n\in R$, with
$a\in\mathbb F_2$ and $n\in N$, one has $Q(r)=0$. Monic division
by $Q$ reduces the value condition to a polynomial $H$ of degree
below $4$. Since
$H(a+n)=H(a)+H'(a)n$, the colon identity says precisely that
\[
                        H(0),H(1),H'(0),H'(1)\in R.
\]
The Chinese remainder theorem for $(X^2)$ and $((X-1)^2)$ identifies
these four coordinates with the four coefficients, over
$\mathbb F_2$ and hence over both coefficient rings. Thus
\begin{equation}\label{eq:forest-module-splitting}
                  \Omega(R,S)=R[X]_{<4}\oplus Q(X)S[X].
\end{equation}

In several variables use the monic basis
$b_i=X^i$ for $i<4$ and $b_i=QX^{i-4}$ for $i\ge4$ in each
coordinate. Terms with a high index vanish on $R^n$ and have
arbitrary $S$-coefficients. For the remaining polynomial, take
finite differences about $a\in\{0,1\}^n$ in a set $J$ of variables
with arbitrary increments $n_j\in N$. The difference equals
$(\partial_JH)(a)\prod_{j\in J}n_j$. Its values lie in $R$;
the products generate $N^{|J|}$, so the colon identity places each
mixed derivative in $R$. The case $J=\varnothing$ is evaluation.
The tensor product of the four-coordinate Hermite isomorphisms
therefore puts all coefficients of the low-index remainder in $R$.

Hence $\Omega_n$ has coefficient module $R$ at every all-low index
and $S$ at every other index. The tensor power of
\eqref{eq:forest-module-splitting} has precisely the same modules,
since $S\otimes_R S=S$. The polynomial tensor map is the identity
on these coefficient modules, proving the asserted isomorphisms.
For reference, the substitutions on $Q$ are
\[
 Q(X+Y)=Q(X)+Q(Y),\qquad
 Q(XY)=Q(X)Q(Y)+Q(X)Y^2+X^2Q(Y).
\]
The all-variable isomorphisms transport every polynomial identity
and composition identity, yielding the unique compatible plethory.

If $S$ were flat, the local inclusion with nonzero residue fibre
would be faithfully flat. A faithfully flat epimorphism is an
isomorphism, contradicting $y_v\notin R$; see
\cite[Lemmas~10.39.16 and 10.107.7]{StacksFaithfulEpimorphisms}.
The split copy of $S$ in \eqref{eq:forest-module-splitting} then
shows that $\Omega(R,S)$ is not flat. Finally $cQ$ is integer-valued
with leading coefficient $c$ for every $c\in S$, so the effective
coefficient ring is all of $S$.
\end{proof}

\begin{proposition}\label{prop:forest-syzygy-invariants}
For the preceding inclusion, put $D=S/R$ and use the modules and
maps of Theorem~\ref{result-20}. Then
\[
                   K_1=K_2=0,\qquad T_1\simeq D^{16}.
\]
The obstruction map records, up to its stipulated sign, the sixteen
low-low coefficients modulo $R$. Both polynomial substitutions have
zero obstruction. No vanishing of $\operatorname{Tor}_q^R(S,S)$
is required for this calculation.
\end{proposition}
\begin{proof}
In the monic basis just used,
\[
 \mathcal F=\bigoplus_{i\ge0}Sb_i,\qquad
 E=\bigoplus_{i<4}Db_i,
\]
and $\varepsilon$ is the quotient $S\to D$ on the four low
summands and zero on the others. Let
$M_q=\operatorname{Tor}_q^R(S,S)$ for $q\ge1$.
The exact sequence $0\to R\to S\to D\to0$ and epicness give
\[
 M_q\simeq\operatorname{Tor}_q^R(S,D)
       \simeq\operatorname{Tor}_q^R(D,S),
\]
as well as
\[
 0\to M_1\to\operatorname{Tor}_1^R(D,D)\to D\to0,
 \qquad \operatorname{Tor}_q^R(D,D)\simeq M_q\quad(q\ge2).
\]
The degree-one row of the universal construction splits by pairs
of basis indices. At a low-low pair it is the diagonal map of
$M_1$, followed by the difference and the displayed injection into
$\operatorname{Tor}_1(D,D)$. Its middle homology is zero and its
cokernel is $D$. At a mixed pair the first map is an isomorphism;
at a high-high pair the other two modules vanish. Hence $K_1=0$
and $T_1=D^{16}$. In degree two the final row map is surjective
at every pair, giving $T_2=K_2=0$.
The cokernel identifies the obstruction with the low-low
coefficients modulo $R$. In the two displayed identities for $Q$,
every term has a high index in at least one coordinate. Multiplying
by the substituted coefficient polynomial preserves this property.
Ordinary $R$-polynomials have low-low coefficients in $R$.
Therefore both substitutions have zero obstruction, as claimed.
\end{proof}
